% Generated draft: requires mathematical and source review.
\documentclass[11pt]{article}
\usepackage[margin=25mm]{geometry}
\usepackage{fontspec}
\setmainfont{Latin Modern Roman}
\usepackage{amsmath,amssymb,mathtools,unicode-math}
\setmathfont{Latin Modern Math}
\usepackage{enumitem,needspace,xurl,hyperref}
\setlength{\parindent}{0pt}
\setlength{\parskip}{5pt}
\setlength{\emergencystretch}{3em}
\newcommand{\N}{\mathbb{N}}
\newcommand{\Z}{\mathbb{Z}}
\newcommand{\Q}{\mathbb{Q}}
\newcommand{\R}{\mathbb{R}}
\newcommand{\C}{\mathbb{C}}
\newcommand{\F}{\mathbb{F}}
\newcommand{\A}{\mathbb{A}}
\newcommand{\PP}{\mathbb{P}}
\newcommand{\E}{\mathbb{E}}
\newcommand{\eps}{\varepsilon}
\newcommand{\dd}{\,\mathrm{d}}
\begin{document}
\section*{UnsolvedMath ID: 20003332}
\textbf{Source problem code:} catalogue entry
\subsection{Definitions and mathematical statement}
Conjecture 41, attributed to J. Marklof, is introduced in the neighboring source paragraph as a prediction for generic negatively curved compact Riemannian manifolds, in the Berry random-wave discussion. It is not an assertion for every compact or finite-volume manifold in the broader spectral setup. The meaning of generic, admissible boundary cases and eigenfunction choices must be fixed. Let $dV$ be normalized volume, $dV(X)=1$, and let $\varphi_i$ be real Laplace eigenfunctions with $\|\varphi_i\|_{L^2(dV)}=1$ and eigenvalues tending to infinity. The conjectured value-distribution limit, for every fixed finite $a<b$, is
\[
 dV\{x:a\leq\varphi_i(x)\leq b\}
 \longrightarrow\frac1{\sqrt{2\pi}}\int_a^b e^{-t^2/2}\,dt.
\]
Real-valuedness is required for these ordered intervals. On a closed manifold positive-eigenvalue functions also have mean zero by orthogonality to constants; that fact is not imposed on unspecified Dirichlet problems.

Let $\mu_i=(\varphi_i)_*dV$ be the probability law of values, and let $\gamma$ be standard Gaussian measure. The statement is exactly $\mu_i\Rightarrow\gamma$. Its characteristic-function target is
\[
 \eps_i(T)=\sup_{|t|\leq T}
 \left|\int_X e^{it\varphi_i(x)}dV(x)-e^{-t^2/2}\right|
                  \longrightarrow0\quad\text{for every fixed }T>0.
\]
The aid states this compact-uniform convergence is equivalent to weak convergence. It is a scalar value-law criterion, not a condition on spatial quantum limits.

The supplied quantitative smoothing lemma asserts that for any probability law $\mu$, $I=[a,b]$ of length $L$, and $\delta,T>0$, with $\eps_\mu(T)=\sup_{|t|\leq T}|\widehat\mu(t)-e^{-t^2/2}|$,
\[
 |\mu(I)-\gamma(I)|\leq
 \frac{2\delta}{\sqrt{2\pi}}
 +\frac{T(L+\delta)}\pi\eps_\mu(T)
 +\frac8{\pi\delta T}.
\]
The terms are Gaussian boundary smoothing, compact-frequency error and high-frequency remainder. The order of choices matters: choose a small smoothing width, then a sufficiently large finite cutoff, then control that cutoff's characteristic-function error. The lemma does not establish this error estimate for the eigenfunctions.

The aid also demonstrates the limit of moment conclusions. For $k\geq2$ put $\epsilon_k=k^{-3}$, $\sigma_k^2=(1-k^{-1})/(1-k^{-3})$, and
\[
 \nu_k=(1-\epsilon_k)N(0,\sigma_k^2)
          +\frac{\epsilon_k}{2}(\delta_k+\delta_{-k}).
\]
These laws have exactly mean zero and variance one, converge weakly to $\gamma$, but have fourth moment $3(1-\epsilon_k)\sigma_k^4+\epsilon_kk^4\to\infty$. They are abstract probability laws, not asserted to arise from eigenfunctions. They show that Gaussian weak convergence plus exact $L^2$ normalization alone cannot yield fourth-moment convergence; extra uniform tail control is required.

Finally the one-point value law is distinct from QUE, which tests the spatial measures $|\varphi_i|^2dV$, and from multi-point random-wave limits, which require joint distributions and correlations. The supplied aid gives these reductions and limitations, not a proof of the conjecture for the intended generic geometric class.
\subsection{Short English statement}
Marklof's conjecture predicts standard Gaussian value distributions for normalized real high-energy eigenfunctions in the intended generic negatively curved compact setting. The supplied aid gives an equivalent characteristic-function target and an explicit smoothing estimate. It also shows why weak Gaussian laws and $L^2$ normalization do not by themselves control fourth moments or establish QUE or multi-point random-wave behavior.
\subsection{Sources}
\begin{itemize}
\item Local UnsolvedMath registry, permanent ID $20003332$, source code catalogue entry; canonical location \texttt{aim-other-notes.json}, \texttt{notes[138]}.
\item AIM \emph{Emerging applications of measure rigidity}, November 1, 2004, Conjecture 41 (J. Marklof), \url{https://aimath.org/WWN/measrigid/measrigid.pdf}. The supplied verification restores the fraction, integration limits and normalized-volume conventions lost in extraction.
\item Neighboring local record $20003331$ supplies the phrase ``generic negatively curved compact Riemannian manifolds'' and M. Berry [19] as random-wave motivation; no full bibliography entry is supplied.
\item Research attempt 1 and summary are marked partial result, proved, and low novelty confidence, with an explicit machine-generated, non-peer-reviewed research-aid caveat requiring independent verification before citation.
\end{itemize}
\end{document}
